% ================================================================
% chapters/chapter2.tex --- RECONSTRUCTION (v2, remarques encadrant)
% ----------------------------------------------------------------
% APPLIQUÉ ICI :
%   - "à/au la plateforme" -> "sur/sur le" (Superadmin, Staff Admin, Company)
%   - "PIN bcrypt" reformulé (clarté)
%   - "Vivre" -> "Lancer" ; "Voir" -> "Consulter" (cohérent avec le
%     reste du tableau qui utilise déjà "Consulter" pour Company/
%     Staff Admin --- À CONFIRMER, l'annotation originale était
%     ambiguë entre les deux mots)
%   - "2.2.3 Besoins non-fonctionnels selon ISO/IEC 25010" ->
%     "2.2.3 Besoins non-fonctionnels" (retrait du qualificatif)
%   - Sous-titres ISO 25010 traduits en français (terme FR d'abord,
%     original anglais entre parenthèses), exactement les termes
%     donnés par l'encadrant
%   - Backlog produit (23 US) : colonne "Titre" ajoutée, chaque User
%     Story reformatée en 3 parties (En tant que / Je veux / Afin de)
%     + critères d'acceptation (DOD), colonne "Épic" retirée du
%     tableau (déjà visible via les séparateurs multicolonnes EP-0X)
%   - "TABLEAU pas table" : gérée globalement dans main.tex
%     (\renewcommand{\tablename}{Tableau}), rien à changer ici
%
% VOLONTAIREMENT PAS AJOUTÉ / EN ATTENTE :
%   - Fonctionnalité "Ajouter un casque" (Company) : question de
%     l'encadrant sur une fonctionnalité potentiellement manquante,
%     pas encore confirmée comme existante côté code --- non ajoutée
%     pour ne pas décrire une fonctionnalité qui n'existe pas
%   - Remarques employé post-simulation : même chose, question
%     ouverte, pas ajoutée
%   - Tableau épics/sprints (2.5.3) : je propose ici une résolution
%     légère (notebox explicatif) plutôt qu'une refonte complète,
%     À VALIDER avec toi
% ================================================================

\chapter{Spécification et pilotage du projet}

\section{Introduction}

Ce chapitre cadre le périmètre fonctionnel et technique de \tynassit{}. Nous identifions
les acteurs, détaillons les besoins fonctionnels et non-fonctionnels selon le modèle
\textbf{ISO/IEC 25010}, présentons la modélisation UML du système, l'architecture retenue,
puis le backlog produit avec priorisation \textbf{MoSCoW} et la planification des sprints.

\section{Analyse des besoins}

\subsection{Identification des acteurs}

Quatre acteurs interagissent avec \tynassit{}~:

\begin{itemize}
  \item \textbf{Superadmin}~: membre de l'équipe Tynass disposant d'un accès total à
    toutes les fonctionnalités, y compris la gestion des comptes administrateurs.
  \item \textbf{Staff Admin}~: membre de l'équipe Tynass chargé de la gestion quotidienne
    des entreprises clientes, des formations et des casques.
  \item \textbf{Company}~: entreprise cliente B2B de Tynass, gérant ses employés et
    consultant les résultats de leurs formations.
  \item \textbf{Employé}~: utilisateur final du casque Meta Quest~3, réalisant les
    simulations de formation.
\end{itemize}

\subsection{Besoins fonctionnels}

Le tableau~\ref{tab:besoins_fonctionnels} décrit les besoins fonctionnels par acteur.

\begin{table}[H]
  \centering
  \caption{Description des besoins fonctionnels par acteur}
  \label{tab:besoins_fonctionnels}
  \begin{tabularx}{\textwidth}{lX}
    \toprule
    \textbf{Acteur} & \textbf{Besoins fonctionnels} \\
    \midrule
    \multirow{4}{*}{\textbf{Superadmin}} &
      S'authentifier sur la plateforme web. \\
    & Gérer les comptes Staff Admin (créer, activer, désactiver). \\
    & Accéder à toutes les fonctionnalités du Staff Admin. \\
    & Créer le premier compte Superadmin via l'endpoint \texttt{/setup}. \\
    \midrule
    \multirow{6}{*}{\textbf{Staff Admin}} &
      S'authentifier sur la plateforme web. \\
    & Gérer les entreprises clientes (CRUD, activation/désactivation). \\
    & Gérer le catalogue de formations VR (CRUD, sceneName). \\
    & Activer un casque Meta Quest~3 via le système Meta User ID. \\
    & Assigner des formations à une entreprise. \\
    & Consulter les statistiques et sessions de toutes les entreprises. \\
    \midrule
    \multirow{4}{*}{\textbf{Company}} &
      S'authentifier sur le dashboard entreprise. \\
    & Gérer les employés (CRUD, génération accessCode, définition d'un
      PIN sécurisé haché en base). \\
    & Révoquer un casque Meta Quest~3. \\
    & Consulter les sessions et performances de ses employés. \\
    \midrule
    \multirow{5}{*}{\textbf{Employé}} &
      Se connecter au casque (accessCode XX-001 + PIN 4 chiffres). \\
    & Accéder aux formations assignées à son entreprise. \\
    & Suivre le tutoriel d'initiation à l'utilisation des contrôleurs
      (scène trailler). \\
    & Lancer la simulation VR (scène intigo). \\
    & Consulter ses résultats envoyés automatiquement à \tynassit{}. \\
    \bottomrule
  \end{tabularx}
\end{table}

\subsection{Besoins non-fonctionnels}

L'analyse des besoins non-fonctionnels s'appuie sur le modèle de qualité logicielle
ISO/IEC~25010~\cite{iso25010}, qui définit huit caractéristiques de qualité.

\placeholder{Modèle ISO/IEC 25010 (figure)}

\subsubsection{Efficience de la performance --- Comportement temporel (Performance
  Efficiency, Time Behaviour)}

Le système doit garantir des temps de réponse acceptables sur toutes ses composantes~:

\begin{itemize}
  \item \textbf{API backend}~: temps de réponse moyen inférieur à \textbf{500~ms}
    pour les endpoints critiques (login, submission de session).
  \item \textbf{Application VR}~: maintien d'une cadence de \textbf{72~fps stables}
    sur Meta Quest~3 (seuil minimum pour éviter le motion sickness).
  \item \textbf{Chargement scènes Unity}~: inférieur à \textbf{5 secondes}.
\end{itemize}

\subsubsection{Sécurité (Security)}

Conformément au principe de défense en profondeur, la sécurité est organisée en
\textbf{quatre couches} complémentaires (détaillées au chapitre~3).

\begin{itemize}
  \item \textbf{Confidentialité}~: JWT signé HS256, différencié par type d'acteur.
  \item \textbf{Résistance}~: bcrypt 10 rounds ($\approx$100~ms), SHA-256 pour les
    tokens sensibles, validation des entrées côté Unity et backend.
\end{itemize}

\subsubsection{Fiabilité (Disponibilité) --- Reliability (Availability)}

\begin{itemize}
  \item \textbf{MongoDB Atlas}~: SLA de 99,9\%, sauvegardes automatiques, accès global.
  \item \textbf{Render + Vercel}~: déploiements continus avec rollback automatique.
\end{itemize}

\subsubsection{Maintenabilité (Modifiabilité) --- Maintainability (Modifiability)}

\begin{itemize}
  \item \textbf{sceneName dynamique}~: l'ajout d'une nouvelle formation VR ne nécessite
    aucune modification du code Unity --- seul le champ \texttt{sceneName} en base de données
    suffit.
  \item \textbf{Architecture MVC}~: modification d'un contrôleur sans impact sur les autres.
  \item \textbf{AppSession statique}~: état global Unity accessible sans instanciation.
\end{itemize}

\subsubsection{Utilisabilité (Apprentissage) --- Usability (Learnability)}

\begin{itemize}
  \item \textbf{Tutoriel VR}~: scène trailler avec 3 leçons couvrant les interactions
    fondamentales (mouvement, manipulation, interface UI).
  \item \textbf{Interface VR simple}~: login en 2 étapes~--- un clavier virtuel
    simplifié et à saisie minimale pour le code d'accès alphanumérique
    (étape~1), puis un pavé numérique dédié, sans clavier virtuel, pour le
    PIN à 4 chiffres (étape~2).
\end{itemize}

\section{Modélisation du système}

\subsection{Vue fonctionnelle --- Diagramme de cas d'utilisation}

\placeholder{Diagramme de cas d'utilisation global (à insérer)}

\subsection{Vue structurelle --- Diagramme de classes}

\placeholder{Diagramme de classes --- 7 modèles MongoDB (à insérer)}

\subsection{Vue dynamique --- Diagrammes de séquence}

\placeholder{Diagramme de séquence --- Activation casque (à insérer)}

\placeholder{Diagramme de séquence --- Login employé VR (à insérer)}

\placeholder{Diagramme de séquence --- SubmitSession (à insérer)}

\section{Architecture du système}

\subsection{Architecture logique}

\tynassit{} adopte une \textbf{architecture en trois couches}~:

\begin{itemize}
  \item \textbf{Couche présentation}~: dashboards React (Admin et Company) et application
    VR Unity~6 sur Meta Quest~3.
  \item \textbf{Couche métier}~: API REST Node.js/Express avec contrôleurs MVC.
  \item \textbf{Couche données}~: MongoDB Atlas avec 7 collections.
\end{itemize}

\placeholder{Architecture logique --- diagramme en couches (à insérer)}

\subsection{Architecture physique --- Diagramme de déploiement}

\placeholder{Diagramme de déploiement (à insérer)}

\noindent Le système est déployé sur trois environnements cloud~:

\begin{itemize}
  \item \textbf{MongoDB Atlas}~: base de données cloud, accessible depuis tous les composants.
  \item \textbf{Render}~: hébergement du backend Node.js avec déploiement continu depuis
    GitHub.
  \item \textbf{Vercel}~: hébergement des deux dashboards React avec CI/CD automatique.
  \item \textbf{Meta Quest~3}~: application Unity déployée en APK via Meta Developer Hub.
\end{itemize}

\section{Pilotage du projet avec Scrum}

\subsection{Application de Scrum}

\placeholder{Processus Scrum adapté au projet (figure)}

\subsection{Backlog produit}

Le backlog produit regroupe \textbf{23 user stories} organisées en \textbf{5 épics},
priorisées selon la méthode \textbf{MoSCoW}~\cite{moscow2009}. Chaque user story suit
le format~: \textit{« En tant que [rôle], je veux [action], afin de [bénéfice]. »},
accompagné de ses critères d'acceptation (DOD --- \textit{Definition of Done}).

\begin{longtable}{lp{2.8cm}p{8cm}l}
  \caption{Backlog produit --- \tynassit{}}
  \label{tab:backlog} \\
  \toprule
  \textbf{ID} & \textbf{Titre} & \textbf{User Story} & \textbf{MoSCoW} \\
  \midrule
  \endfirsthead
  \toprule
  \textbf{ID} & \textbf{Titre} & \textbf{User Story} & \textbf{MoSCoW} \\
  \midrule
  \endhead
  \bottomrule
  \endfoot

  \multicolumn{4}{l}{\textit{EP-01 --- Backend \& Sécurité}} \\[2pt]

  US-01 & Authentification sécurisée (admin) &
    En tant qu'admin, je veux m'authentifier sur la plateforme, afin
    d'accéder aux fonctionnalités de gestion réservées à mon rôle.
    \newline\textbf{Critères d'acceptation}~:
    \begin{itemize}
      \item Le système vérifie l'email et le mot de passe.
      \item Un message d'erreur s'affiche en cas d'identifiants incorrects.
      \item Un token JWT est généré au succès.
      \item Le mot de passe est haché en base (bcrypt).
    \end{itemize}
    & \must{} \\[4pt]

  US-02 & Gestion des entreprises clientes &
    En tant qu'admin, je veux créer et gérer les entreprises clientes,
    afin de leur donner accès à la plateforme.
    \newline\textbf{Critères d'acceptation}~:
    \begin{itemize}
      \item CRUD complet (création, modification, désactivation).
      \item Le nom et l'email de l'entreprise sont obligatoires et validés.
      \item Une entreprise désactivée ne peut plus se connecter.
    \end{itemize}
    & \must{} \\[4pt]

  US-03 & Gestion du catalogue de formations &
    En tant qu'admin, je veux créer des formations VR avec leur scène
    associée, afin de les rendre disponibles aux entreprises clientes.
    \newline\textbf{Critères d'acceptation}~:
    \begin{itemize}
      \item Le champ \texttt{sceneName} est obligatoire à la création.
      \item La formation est visible dans le catalogue dès sa création.
      \item Sa modification ne nécessite aucun redéploiement Unity.
    \end{itemize}
    & \must{} \\[4pt]

  US-04 & Activation d'un casque via Meta User ID &
    En tant qu'admin, je veux activer un casque Meta Quest~3 via un
    code généré par le système, afin de sécuriser l'accès physique à
    la plateforme.
    \newline\textbf{Critères d'acceptation}~:
    \begin{itemize}
      \item Un code à 6 chiffres est généré et expire après 15 minutes.
      \item Un token device est généré au succès et haché SHA-256 en base.
      \item Le casque est associé à une entreprise précise.
    \end{itemize}
    & \must{} \\[4pt]

  US-05 & Consultation des statistiques globales &
    En tant qu'admin, je veux consulter les statistiques globales de
    la plateforme, afin de suivre l'usage et la performance des
    formations.
    \newline\textbf{Critères d'acceptation}~:
    \begin{itemize}
      \item Les données sont agrégées par formation et par entreprise.
      \item Les statistiques sont recalculées à chaque requête.
      \item L'accès est réservé au rôle admin.
    \end{itemize}
    & \should{} \\[4pt]

  \multicolumn{4}{l}{\textit{EP-02 --- Dashboard Admin}} \\[2pt]

  US-06 & Connexion au dashboard Admin &
    En tant qu'admin, je veux me connecter au dashboard web, afin
    d'accéder à l'interface de gestion.
    \newline\textbf{Critères d'acceptation}~:
    \begin{itemize}
      \item Formulaire email et mot de passe.
      \item Le token est conservé côté client entre les pages.
      \item Redirection automatique vers la connexion si la session expire.
    \end{itemize}
    & \must{} \\[4pt]

  US-07 & Visualisation des KPIs &
    En tant qu'admin, je veux voir les indicateurs clés (entreprises,
    sessions, taux de réussite) sur mon tableau de bord, afin d'avoir
    une vue d'ensemble de l'activité.
    \newline\textbf{Critères d'acceptation}~:
    \begin{itemize}
      \item 4 KPIs sont affichés à l'ouverture du dashboard.
      \item Le top 5 des entreprises et formations les plus actives
        est visible.
      \item La page se charge en moins de 2 secondes.
    \end{itemize}
    & \must{} \\[4pt]

  US-08 & Gestion des entreprises depuis le dashboard &
    En tant qu'admin, je veux gérer les entreprises et leur assigner
    des formations depuis le dashboard, afin de configurer l'accès de
    chaque client.
    \newline\textbf{Critères d'acceptation}~:
    \begin{itemize}
      \item CRUD complet accessible depuis l'interface.
      \item L'assignation de formation se fait via un dialog dédié.
      \item Le changement est visible immédiatement côté entreprise.
    \end{itemize}
    & \must{} \\[4pt]

  US-09 & Activation et révocation des casques &
    En tant qu'admin, je veux activer ou révoquer un casque depuis le
    dashboard, afin de contrôler les appareils actifs sur la
    plateforme.
    \newline\textbf{Critères d'acceptation}~:
    \begin{itemize}
      \item L'activation se fait via le code à 6 chiffres affiché sur
        le casque.
      \item La révocation invalide le token immédiatement.
      \item Un casque révoqué ne peut plus s'authentifier.
    \end{itemize}
    & \must{} \\[4pt]

  US-10 & Filtrage des sessions &
    En tant qu'admin, je veux filtrer les sessions par entreprise,
    formation ou date, afin de retrouver rapidement une information
    précise.
    \newline\textbf{Critères d'acceptation}~:
    \begin{itemize}
      \item Les filtres sont combinables entre eux.
      \item Les résultats se mettent à jour sans rechargement de page.
      \item Les filtres peuvent être réinitialisés en un clic.
    \end{itemize}
    & \should{} \\[4pt]

  \multicolumn{4}{l}{\textit{EP-03 --- Dashboard Company}} \\[2pt]

  US-11 & Connexion au dashboard Company &
    En tant qu'entreprise, je veux me connecter à mon dashboard, afin
    de gérer mes employés et consulter mes résultats.
    \newline\textbf{Critères d'acceptation}~:
    \begin{itemize}
      \item Authentification par email et mot de passe propres à
        l'entreprise.
      \item L'accès est limité aux seules données de cette entreprise.
    \end{itemize}
    & \must{} \\[4pt]

  US-12 & Gestion des employés &
    En tant qu'entreprise, je veux gérer mes employés (créer,
    modifier, définir un PIN), afin de leur donner accès aux
    formations sur le casque.
    \newline\textbf{Critères d'acceptation}~:
    \begin{itemize}
      \item L'accessCode est généré automatiquement (format XX-001).
      \item Le PIN saisi est haché avant stockage en base.
      \item Un employé désactivé ne peut plus se connecter au casque.
    \end{itemize}
    & \must{} \\[4pt]

  US-13 & Consultation des sessions employés &
    En tant qu'entreprise, je veux consulter les sessions de mes
    employés, afin de suivre leur progression et leurs résultats.
    \newline\textbf{Critères d'acceptation}~:
    \begin{itemize}
      \item Liste paginée des sessions.
      \item Filtrage par employé ou par formation.
      \item Score et statut (réussi/échoué) visibles pour chaque session.
    \end{itemize}
    & \should{} \\[4pt]

  \multicolumn{4}{l}{\textit{EP-04 --- Application VR Unity}} \\[2pt]

  US-14 & Tutoriel des contrôleurs VR &
    En tant qu'employé, je veux apprendre à utiliser les contrôleurs
    VR via un tutoriel, afin d'être à l'aise avant ma première
    formation.
    \newline\textbf{Critères d'acceptation}~:
    \begin{itemize}
      \item 3 leçons couvrant déplacement, manipulation et interaction UI.
      \item Possibilité de passer le tutoriel.
      \item Transition automatique vers l'écran de connexion à la fin.
    \end{itemize}
    & \must{} \\[4pt]

  US-15 & Activation du casque &
    En tant qu'employé, je veux que le casque s'active automatiquement
    via son identifiant matériel, afin de ne pas ressaisir un code à
    chaque session.
    \newline\textbf{Critères d'acceptation}~:
    \begin{itemize}
      \item Le Meta User ID est envoyé au backend au démarrage.
      \item Un casque déjà activé est reconnu directement.
      \item Un casque non activé affiche l'écran d'attente.
    \end{itemize}
    & \must{} \\[4pt]

  US-16 & Connexion employé sur le casque &
    En tant qu'employé, je veux me connecter avec mon code d'accès et
    mon PIN, afin d'accéder à mes formations de façon sécurisée.
    \newline\textbf{Critères d'acceptation}~:
    \begin{itemize}
      \item Le code est validé au format XX-001 en temps réel.
      \item Le PIN à 4 chiffres est masqué à la saisie.
      \item Un message d'erreur s'affiche en cas d'identifiants incorrects.
    \end{itemize}
    & \must{} \\[4pt]

  US-17 & Consultation des formations disponibles &
    En tant qu'employé, je veux voir la liste des formations assignées
    à mon entreprise, afin de choisir celle à suivre.
    \newline\textbf{Critères d'acceptation}~:
    \begin{itemize}
      \item Les formations sont affichées en grille.
      \item Seules les formations assignées à l'entreprise sont visibles.
      \item La sélection lance directement la scène correspondante.
    \end{itemize}
    & \must{} \\[4pt]

  US-18 & Configuration des paramètres &
    En tant qu'employé, je veux configurer l'audio, la langue et les
    contrôleurs, afin d'adapter l'expérience à mes préférences.
    \newline\textbf{Critères d'acceptation}~:
    \begin{itemize}
      \item Les réglages sont sauvegardés localement.
      \item Ils sont appliqués immédiatement.
      \item Ils sont conservés entre les sessions.
    \end{itemize}
    & \should{} \\[4pt]

  \multicolumn{4}{l}{\textit{EP-05 --- Simulation VR Indigo}} \\[2pt]

  US-19 & Exploration des locaux &
    En tant qu'employé, je veux explorer virtuellement les locaux
    d'Indigo et ses différents départements, afin de me familiariser
    avec l'environnement avant ma prise de poste.
    \newline\textbf{Critères d'acceptation}~:
    \begin{itemize}
      \item Les 7 départements sont accessibles dans la scène.
      \item Un panneau d'information est disponible à chaque zone.
      \item Une narration audio accompagne chaque département.
    \end{itemize}
    & \must{} \\[4pt]

  US-20 & Quiz de fin de simulation &
    En tant qu'employé, je veux répondre à un quiz sur ce que j'ai
    appris pendant l'exploration, afin de vérifier ma compréhension.
    \newline\textbf{Critères d'acceptation}~:
    \begin{itemize}
      \item Une question par département visité.
      \item Le score est calculé en pourcentage.
      \item Le seuil de réussite est fixé à 70\%.
    \end{itemize}
    & \must{} \\[4pt]

  US-21 & Soumission automatique des résultats &
    En tant que système, je veux soumettre automatiquement les
    résultats de la simulation au backend, afin que l'entreprise
    puisse suivre la progression de l'employé.
    \newline\textbf{Critères d'acceptation}~:
    \begin{itemize}
      \item L'envoi est déclenché automatiquement à la fin du quiz.
      \item La session est liée à l'employé et à la formation.
      \item Le numéro de tentative est incrémenté automatiquement.
    \end{itemize}
    & \must{} \\[4pt]

  US-22 & Affichage des résultats &
    En tant qu'employé, je veux voir mon score et le détail de mes
    résultats à la fin de la simulation, afin de savoir si j'ai
    réussi.
    \newline\textbf{Critères d'acceptation}~:
    \begin{itemize}
      \item Le score est affiché en pourcentage.
      \item Le statut réussi/échoué est visible.
      \item Un retour vers la liste des formations est possible.
    \end{itemize}
    & \must{} \\[4pt]

  US-23 & Navigation dans la liste des formations &
    En tant qu'employé, je veux faire défiler la liste des formations
    avec le joystick, afin de naviguer confortablement quand il y en
    a beaucoup.
    \newline\textbf{Critères d'acceptation}~:
    \begin{itemize}
      \item Le défilement au joystick est fluide.
      \item Il fonctionne avec plus de 10 formations affichées.
      \item Il n'entre pas en conflit avec la sélection tactile.
    \end{itemize}
    & \could{} \\
\end{longtable}

\subsection{Planification des sprints}

\begin{table}[H]
  \centering
  \caption{Planification des sprints}
  \label{tab:sprint_plan}
  \begin{tabular}{llll}
    \toprule
    \textbf{Sprint} & \textbf{Titre} & \textbf{Durée} & \textbf{US couvertes} \\
    \midrule
    Sprint 1 & Backend Core \& Sécurité          & 3 semaines & US-01 à US-05 \\
    Sprint 2 & Dashboards Admin \& Company        & 3 semaines & US-06 à US-13 \\
    Sprint 3 & Application VR Unity (Auth)        & 4 semaines & US-14 à US-18 \\
    Sprint 4 & Simulation VR Indigo (intigo)      & 5 semaines & US-19 à US-23 \\
    \bottomrule
  \end{tabular}
\end{table}

\begin{notebox}
  La durée d'un sprint (3, 3, 4 et 5 semaines) a été fixée au début du projet et
  maintenue constante dans son principe de \textit{timeboxing}~: chaque sprint
  regroupe les user stories d'un même épic parce qu'elles présentent une forte
  dépendance technique séquentielle (le dashboard ne peut pas être développé avant
  que le backend existe, par exemple), et non parce qu'un épic correspondrait par
  définition à un seul sprint --- ce n'est pas le cas en général, et cela reste ici
  une simplification propre à la petite échelle de ce projet (un seul développeur,
  quatre sprints).
\end{notebox}

\section{Conclusion}

Ce chapitre a établi les bases de \tynassit{}~: identification des acteurs, besoins
fonctionnels et non-fonctionnels (structurés selon ISO/IEC~25010), modélisation UML,
architecture déployée et backlog MoSCoW de 23 user stories. Le chapitre suivant présente
la réalisation du premier sprint~: le backend core et son architecture de sécurité.
